% Ad Familiares 7.24 --- headnote
% ad-familiares-07-24, c. 20 August 45 BC, Tusculum

\textit{Cicero to M. Fadius Gallus, written from the Tusculan villa
around 20 August 45 BC. The matter at hand is the singer Tigellius
Sardus --- one of Caesar's intimates, lampooned later by Horace ---
whose ill-will toward Cicero, Gallus had been working (a touch
fussily, in Cicero's view) to repair. Cicero is having none of it.
He cites Cipius, the legendary Roman husband who pretended to be
asleep at the right moments: ``I am not asleep for everyone''; just
so, ``I am not at every man's service.'' The pose recovered: in
his consular pomp he was less besieged by suitors than he now is,
in retirement, by Caesarians --- Tigellius alone excepted, whom he
will count as so much clear gain forgone, and whom he supposes Calvus
Licinius (the iambic poet, in the manner of Hipponax) already has
fully bought up in scurrilous verse.}

\textit{The second section gives the specific grievance. Cicero
had taken on the case of his friend Phamea, but the trial day
Phamea proposed collided with the panel for P. Sestius; Cicero
offered any other day Phamea liked. Phamea --- relying, as Cicero
notes acidly, on a grandson who was a passably good piper and a
masseur of his own --- went off in a huff. So much, Cicero says,
for ``Sardinians for sale, each worse than the next,'' the proverbial
tag (the slave-market joke about Sardus the captive having sired
Sardinian Tigellius). The Perseus dateline is \textit{circa xiii
K. Sept.}, i.e.\ about 20 August 45 BC. The letter closes by
asking Gallus for his \textit{Cato} --- Gallus had written one of
the encomia of Cato Uticensis that the year's literary contest
prompted --- which Cicero ``ought to be ashamed,'' for both their
sakes, not yet to have read.}
